\chapter{Preterm Labor and Preterm Birth}
\index{Preterm Labor and Preterm Birth}
\label{sec:Preterm Labor and Preterm Birth}

\section{Overview}
Preterm labor is defined as regular uterine contractions causing cervical change between 20 and 36 weeks 6 days of gestation. Preterm birth (delivery before 37 weeks) affects 10\% of pregnancies worldwide and is the leading cause of neonatal mortality and morbidity. The Disease mechanism involves multiple pathways: infection/inflammation (40\% of early preterm births), uteroplacental reduced blood flow, uterine overdistension, and immunologically mediated rejection. This pregnancy-related condition requires careful monitoring to ensure the safety of both mother and baby. Early detection and management are essential.
\section{Causes}
\section{Risk Factors}

\begin{itemize}[noitemsep]
  \item Genetic predisposition and family history
  \item Advancing age
  \item Lifestyle factors (diet, exercise, smoking, alcohol)
  \item Environmental exposures
  \item Underlying medical conditions
  \item Medication use
\end{itemize}
\section{Signs and Symptoms}

\subsection{Common symptoms}
\begin{itemize}[noitemsep]
  \item \item You may experience regular contractions with cervical change before 37 weeks. Pain or discomfort in the affected area    \item Pain or discomfort in the affected area
  \end{itemize}

\subsection{Emergency warning signs}
\begin{itemize}[noitemsep]
  \item Severe or worsening symptoms of preterm labor and preterm birth require immediate medical evaluation
\end{itemize}
\section{Progression and Stages}
The condition may progress through different stages of severity over time. Early stages often involve mild or subtle symptoms that may be overlooked. Moderate stages involve more noticeable symptoms affecting daily function. Advanced stages may involve significant impairment and complications.
\section{Complications}
\begin{itemize}[noitemsep]
  \item Disease progression leading to worsening symptoms
  \item Secondary infections or injuries
  \item Side effects of treatment
  \item Reduced quality of life
  \item Emotional and psychological impact
\end{itemize}
\section{Diagnosis}
\begin{itemize}[noitemsep]
  \item To make the diagnosis, doctors look for documented cervical change on exam
  \item Transvaginal ultrasound measurement of cervical length (<25 mm at 20--24 weeks is high risk) and fetal fibronectin testing help stratify risk. Acute Treatment options include antenatal corticosteroids for fetal lung maturity (24--34 weeks), magnesium sulfate for fetal neuroprotection (24--32 weeks), and tocolysis to delay delivery for 48 hours to allow corticosteroid benefit.
  \item Comprehensive medical history and physical examination
  \item Laboratory tests to assess organ function
  \item Imaging studies to visualize affected structures
  \item Specialized testing depending on the affected system
  \item Consultation with relevant specialists
\end{itemize}
\section{Prevention}
\begin{itemize}[noitemsep]
  \item Maintain a healthy lifestyle with balanced diet and exercise
  \item Avoid known risk factors
  \item Attend regular medical check-ups for early detection
  \item Follow recommended screening guidelines
\end{itemize}
\section{Treatment Options}
Medications to manage symptoms and address underlying mechanisms Lifestyle modifications and self-care measures Physical therapy and rehabilitation Surgical intervention when indicated Regular monitoring and follow-up care Patient education and support services

\begin{itemize}[noitemsep]
  \item Medications to manage symptoms and address underlying mechanisms
  \item Lifestyle modifications and self-care measures
  \item Physical therapy and rehabilitation
  \item Surgical intervention when indicated
  \item Regular monitoring and follow-up care
\end{itemize}
\section{Living with It}
\begin{itemize}[noitemsep]
  \item Follow your treatment plan as prescribed
  \item Maintain regular follow-up with your healthcare provider
  \item Make necessary lifestyle adjustments
  \item Seek support from family, friends, or support groups
  \item Stay informed about your condition
\end{itemize}
\section{Outlook}
Your outlook depends on gestational age at delivery. The outlook varies depending on the specific condition, severity at diagnosis, and response to treatment. Early intervention generally leads to better outcomes. Many conditions can be effectively managed with appropriate treatment. Regular follow-up care is important for monitoring and treatment adjustment.
